\section{Methods}\label{s:setting}

\subsection{Threat model}\label{s:threat}
The verifier owns the probe. From the GPU it receives only the probe's trace, and from
the operator only the declaration. It also holds recordings of the declared work, made
during runs that it observes or controls.
The operator controls the software on the GPU, cannot touch the probe, and may hide
training by the strategies of \cref{s:corpus}. We do not test an operator who adapts to
the detector.

\subsection{Channel}
A single \GPU{} runs every workload. A verifier-owned Rogowski current probe (\Probe{},
\ProbeSens{}~mV/A, \ProbeLo{}~Hz to \ProbeHi{}~MHz) is clamped to the GPU's auxiliary supply. We
sample at \Fs{}~S/s with a \DigBits-bit \Digitiser{} and no anti-alias filter.
\Cref{a:impl} bounds the variance that folds into the band. The probe is AC-coupled. It has no DC response, and there is no
voltage channel. We therefore measure current, not power.

\subsection{Corpus}\label{s:corpus}
We use \NCaptures{} captures with \CorpusHours{}~h of current from a public
corpus \citep{kimpson2026corpus}, all from one
language model: a decoder-only transformer of \ModelLayers{} layers and width
\ModelWidth, with \ModelParams{}M parameters, built directly in PyTorch. Captures fall into workload groups:
\begin{itemize}
  \item \emph{honest training}: full fine-tuning and LoRA. Some full fine-tuning runs add
        checkpoint evaluation or a copying control, or discard every update;
  \item \emph{time-sharing}: training and inference alternate in slots;
  \item \emph{co-location}: training and inference run as two resident processes,
        and the driver time-slices them. Each run sets a target training share
        $\rho \in \{\RhoLo, \RhoMid, \RhoHi\}$, one run each. At $\rho = \RhoLo$
        and $\RhoHi$ one process pauses between steps to hold the share; we call
        these runs \emph{throttled}. At
        $\rho = \RhoMid$ neither process pauses, and the driver alone sets the share.
        We call this run \emph{unthrottled};
  \item \emph{jitter}, \emph{padding} and \emph{power cap}: training with
        perturbed step timing or a lowered power limit;
  \item \emph{declared work}, the workloads of the declaration: serving decode
        (generating tokens), serving prefill (reading prompts) and idle;
  \item \emph{non-AI work}: dense GEMM, memory-bound kernels and periodic decoys.
\end{itemize}
Time-sharing, co-location, jitter, padding and power cap are the \emph{evasion groups}.
We split the captures by whole days, so that no run spans two splits.
$\mathcal{D}_{\mathrm{dev}}$ holds every day before the test day,
$\mathcal{D}_{\mathrm{test}}$ the test day, and $\mathcal{D}_{\mathrm{cal}}$ the
calibration day after it. Co-location exists only in
$\mathcal{D}_{\mathrm{dev}}$. \Cref{t:corpus} counts runs, hours and windows per
workload and split.

Earlier studies use short recordings: \citet{gargiulo2026} record \GkTraceS{}-s traces
at about \GkRateApprox{}~MS/s, one workload each, and \citet{rahman2026} use NVML
windows of \RtWinLo{} to \RtWinHi{}~s. We record continuous hours at \Fs{}~S/s. A
verifier watches a GPU for hours, and an evasion can start and stop inside a recording
of declared work, so only continuous recordings let us test detectors as monitors
(\cref{s:works}).

\subsection{Windows and features}
A window $x$ is \Win{}~s of trace. The hop between consecutive windows is \HopSnap{}~s
for window-level results and \HopMon{}~s for the monitor. We define a window as
\emph{training} iff it holds at least one
complete optimizer step, following \citet{gargiulo2026}. Windows with only part
of a step are left out of detection rates. Each window gives \NFeat{} amplitude-free
features, adapted from \citet{gargiulo2026}, of the raw spectrum and of the envelope
$e(t)$, the RMS of the raw trace over \EnvFrame{}~ms frames.

\subsection{Detectors}\label{s:detectors}
Each detector assigns a score $s(x)$ to a window $x$. Our detectors follow the two
approaches of \citet{scher2025}.

The \emph{signature check} $s_R$ is an isolation
forest~\citep{liu2008}. It is fitted to the declared windows of
$\mathcal{D}_{\mathrm{dev}}$ only: decode, prefill and idle. Its score is the negated isolation score, so a high $s_R$ means ``unlike the
declaration''. It never sees training, evasion or non-AI work.

The \emph{classifier} is the two-stage random forest of \citet{gargiulo2026}, with
\ForestTrees{} trees per stage (settings in \cref{a:impl}). It learns from labelled examples of
training, so it stands for the default approach, even where it is fitted to this GPU and
its declared work. Stage one separates AI from non-AI
work. Stage two then separates training from inference in the windows that stage one
calls AI.
The classifier's score is
\begin{equation}
  s_C(x) = \begin{cases}
    \hat P(\text{training} \mid x, \text{AI}) & \text{stage one calls } x \text{ AI},\\
    0 & \text{otherwise}.
  \end{cases}
\end{equation}
We fit the classifier in three ways:
\begin{itemize}
  \item $s_C^{\mathrm{dev}}$ is fitted to $\mathcal{D}_{\mathrm{dev}}$ only. To test each
        evasion as unseen, we hold its group out of the fit, so there is one fit per
        evasion group; each group is scored by the fit that holds it out.
        ``The classifier'' means $s_C^{\mathrm{dev}}$. Its window-level $\tau_\alpha$
        and false-positive rates come from the fit that holds out co-location, and the
        monitor uses the fit that holds out time-sharing (\cref{s:monitor});
  \item the \emph{no-evasion fit} is $s_C^{\mathrm{dev}}$ with every evasion group
        removed;
  \item the \emph{replica} is fitted to \NRepTraces{} traces
        recorded following the design of \citet{gargiulo2026}, where it reaches
        macro-$F_1 = \SoneFone$ under leave-one-group-out.
\end{itemize}
$s_C^{\mathrm{dev}}$ and the no-evasion fit are fitted to our corpus, so the declared
work is among their examples of work that is not training. The replica is fitted to
none of the \NCaptures{} captures of \cref{s:corpus}: its traces, one workload each, were recorded separately on the
same GPU and probe.

\subsection{Thresholds}
A detector flags $x$ iff $s(x) > \tau_\alpha$, with
\begin{equation}\label{e:tau}
  \tau_\alpha = Q_{1-\alpha}\bigl\{ s(x) : x \in \mathcal{D}_{\mathrm{cal}}^{\mathrm{decl}} \bigr\}, \quad \alpha = \Alpha.
\end{equation}
$Q_{1-\alpha}$ is the empirical quantile, with linear interpolation, and
$\mathcal{D}_{\mathrm{cal}}^{\mathrm{decl}}$ is the declared windows of
$\mathcal{D}_{\mathrm{cal}}$, one decode run and one idle run of equal length. $s_R$ and every
classifier fit use the same $\tau_\alpha$ rule. For a set $W$ of
training windows, detection is
$d = |W|^{-1} \sum_{x \in W} \mathbb{I}[s(x) > \tau_\alpha]$. The false-positive rate
is $d$ with $W$ the declared windows of $\mathcal{D}_{\mathrm{test}}$.

For $s_R$, $\tau_\alpha$ is our choice of rule for what \citet{scher2025} call
``sufficiently close to the baseline''. They
specify one baseline, whereas we use two separate recordings of declared work: the fit
($\mathcal{D}_{\mathrm{dev}}$) and the threshold set
($\mathcal{D}_{\mathrm{cal}}^{\mathrm{decl}}$). Together they form the \emph{reference} of
$s_R$. We tested no other rule. We varied only $\alpha$ (\cref{s:works}).

\subsection{Monitor}\label{s:monitor}
We run each detector over the test day as a monitor (\cref{f:f8}). This represents a verifier that
observes the current trace continuously. The verifier does not raise an alarm on a
single flagged window. Instead, an alarm starts when 3 of 5 consecutive windows are
flagged, and ends after \Rearm{} unflagged windows in a row;
\cref{s:works} varies this rule. We describe false
alarms by two numbers: $F$, the alarm onsets per monitored hour, and $A$, the fraction
of time in alarm.
